\subsection{EPR Paradox and Quantum Correlation}

\subsubsection{Definition of EPR Paradox}

In 1935, Albert Einstein, Boris Podolsky, and Nathan Rosen published a paper entitled \emph{Can Quantum-Mechanical Description of Physical Reality Be Considered Complete?}.\cite{einstein1935can} Their goal was \underline{not} to disprove quantum mechanics, but to \textbf{show} that the \textbf{theory appeared to be incomplete}.

\highspace
The source of their concern was a peculiar phenomenon predicted by quantum mechanics: two particles could become so strongly correlated that measuring one particle would instantaneously determine the state of the other, regardless of the distance separating them (explained in the proof of the no-signaling principle, \autopageref{sec:no-signaling-principle}).

\highspace
Einstein found this deeply troubling. According to special relativity, \textbf{no information can travel faster than the speed of light}. Yet entanglement seemed to imply an instantaneous connection between distant particles, a phenomenon Einstein famously called \emph{``spooky action at a distance''} (\emph{spukhafte Fernwirkung}).

\highspace
The EPR paper sparked one of the most important debates in the history of physics. Shortly after its publication, Erwin Schrödinger analyzed the EPR argument and realized that it revealed what he considered the most distinctive feature of quantum mechanics. In a series of papers published the same year, he introduced the term \emph{Verschränkung} (``entanglement''), which is now known as quantum entanglement (\autopageref{sec:entanglement}).

\highspace
Interestingly, Schrödinger largely shared Einstein's discomfort with the standard interpretation of quantum mechanics. However, while Einstein viewed the EPR argument as evidence that the theory was incomplete, \textbf{Schrödinger recognized entanglement as a fundamental property of quantum systems}.

\highspace
The central \hl{question raised by the EPR paradox} is therefore:
\begin{quote}
    \emph{If measuring Alice's qubit immediately determines the state of Bob's qubit, \textbf{does quantum mechanics allow faster-than-light communication?}}
\end{quote}
As we have already seen through the no-signaling principle (\autopageref{sec:no-signaling-principle}), the \textbf{answer is no}. Entanglement creates correlations between distant measurements, but these correlations cannot be used to transmit information.

\highspace
Precisely, Einstein's concern was fundamentally:
\begin{quote}
    \emph{If measuring Alice's qubit instantaneously determines the state of Bob's qubit, what is actually happening?}
\end{quote}
In other words, \textbf{how can nature produce these instantaneous correlations?} Einstein was not primarily asking about the possibility of faster-than-light communication, but rather about the underlying mechanism that could explain the observed correlations without violating relativity. He suspected that Bob's particle already possessed some hidden information before the measurement, which would determine the outcome of Alice's measurement without any need for instantaneous influence.

\highspace
\begin{flushleft}
    \textcolor{Green3}{\faIcon{book} \textbf{Modern Perception of the EPR Paradox}}
\end{flushleft}
To understand the paradox in a modern quantum-computing setting, consider the Bell state:
\begin{equation*}
    \ket{\Phi^+}
    =
    \frac{\ket{00}+\ket{11}}{\sqrt2}
\end{equation*}
Suppose Alice and Bob each possess one qubit of this entangled pair and are separated by a very large distance.

\highspace
If Alice measures her qubit in the computational basis and obtains:
\begin{equation*}
    \ket{0}
\end{equation*}
Then Bob's qubit immediately collapses to:
\begin{equation*}
    \ket{0}
\end{equation*}
Similarly, if Alice measures:
\begin{equation*}
    \ket{1}
\end{equation*}
Then Bob's qubit immediately collapses to:
\begin{equation*}
    \ket{1}
\end{equation*}
At first sight, this appears to imply an instantaneous influence between distant particles. \hl{How can Bob's qubit ``know'' which outcome Alice obtained if they may be separated by billions of light-years?} This is the modern formulation of the EPR paradox.

\highspace
\begin{definitionbox}[: Einstein-Podolsky-Rosen (EPR) Paradox]
    The \definition{Einstein-Podolsky-Rosen (EPR) Paradox} is the \textbf{apparent contradiction} between quantum \textbf{entanglement} and the \textbf{classical principle of locality}.

    \highspace
    More precisely, if two particles are prepared in an entangled state, a measurement performed on one particle appears to instantaneously determine the state of the other, regardless of the distance separating them.

    \highspace
    This raises the question of whether \hl{quantum mechanics is a complete description of physical reality}, or whether \hl{additional hidden information is required to explain these correlations without instantaneous influence}.
\end{definitionbox}